\documentclass[11pt,a4paper]{letter}
\usepackage[T1]{fontenc}
\usepackage[margin=2.5cm]{geometry}
\usepackage{hyperref}
\signature{Diego Andres Perez Ruiz}
\address{Department of Social Statistics\\School of Social Sciences\\University of Manchester\\Manchester, United Kingdom\\\href{mailto:diego.perezruiz@manchester.ac.uk}{diego.perezruiz@manchester.ac.uk}}
\begin{document}
\begin{letter}{The Editor-in-Chief\\The R Journal}
\opening{Dear Editor,}

I submit the article ``ordpp: Projection Pursuit with Monotone Scoring for Ordinal Data in R'' for consideration as a contributed research article in \emph{The R Journal}.

Ordered categorical items are central to survey research and psychometrics, yet they are usually analysed with equally spaced numerical codes. The article presents the R package \textbf{ordpp}, which treats the numerical spacing of categories as part of an exploratory analysis. It jointly estimates monotone category scores and a projection direction by maximizing a characteristic-function discrepancy from item independence with the observed marginals. The package then freezes the resulting measurement map and evaluates it on independent data, including a permutation test that is valid in finite samples and requires no refitting. A wrapper allows score maps produced by other packages, such as \textbf{Gifi}, \textbf{ordPens} and \textbf{psych}, to be evaluated with exactly the same tools.

I believe the article suits the journal for three reasons. First, it documents a complete, openly available R workflow (fit, inspect, validate, test, check stability, predict) rather than a method alone. Second, it relates the package carefully to existing R software for optimal scaling, latent-response models, projection pursuit and distance covariance, and compares it with several of these packages on equal terms. Third, it reports its limitations candidly. In simulations, principal-component directions gave a more powerful held-out test under weak dependence than the pursued directions, and a larger discrepancy did not imply better recovery of a latent factor. The article therefore positions \textbf{ordpp} as a descriptive tool and shows users how to test along other directions when detection is the goal.

The article includes three short propositions (characterization of independence, a low-frequency expansion of the index, and finite-sample validity of the held-out test), an application to the openly available SAPA personality data in \textbf{psychTools}, and a simulation study of size, power and latent recovery with Monte Carlo standard errors.

The package (version 0.1.0) passes \texttt{R CMD check -{}-as-cran} on macOS and on win-builder (R-devel) with only the ``new submission'' note, and it has been submitted to CRAN. The source code and replication materials are available at \url{https://github.com/dapr12/ordpp} and archived at \url{https://doi.org/10.5281/zenodo.23186811}. The script \texttt{scripts/reproduce\_article.R} regenerates every number, table and figure in the article in a few minutes from the saved outputs of the long-running simulations; the complete computations can be re-run with \texttt{FULL=1}. One comparator, \textbf{ordPens}, had been archived from CRAN at the time of writing and was installed from the CRAN archive, as noted in the article.

The manuscript is original, has not been published elsewhere and is not under consideration by another journal.

\closing{Yours sincerely,}
\end{letter}
\end{document}
